\subsection{INCREASING MAPPINGS}

DEFINITION 1. Let E and F be preordered sets (the preorder relation in each being denoted by \(\leqslant\) ). A mapping f of E into F is said to be increasing (or order-preserving) if the relation \(x \leqslant y\) implies \(f(x) \leqslant f(y)\) ; it is decreasing (or order-reversing) if the relation \(x \leqslant y\) implies \(f(x) \geqslant f(y)\) . A mapping of E into F is said to be monotone if it is either increasing or decreasing.

An increasing mapping of E into F becomes decreasing (and vice versa) when we replace one of the preorderings on E or on F by the opposite preordering. Every constant function is both increasing and decreasing; the converse of this statement is usually not true.

For example, if a set E is ordered by the equality relation, the identity mapping of E onto itself is both increasing and decreasing, but not constant if E has more than one element (cf.~Exercise 7).

DEFINITION 2. Let E and F be two ordered sets. A mapping f of E into F is said to be strictly increasing (or strictly order-preserving) if the relation x \textless{} y implies \(f(x) < f(y)\) ; f is said to be strictly decreasing (or strictly order-reversing) if the relation x \textless{} y implies \(f(x) > f(y)\) . A mapping of E into F is said to be strictly monotone if it is either strictly increasing or strictly decreasing.

Examples

\begin{enumerate}
\def\labelenumi{(\arabic{enumi})}
\item
  Let E be a set. The mapping \(X \rightarrow E - X\) of \(\mathfrak{P}(E)\) (ordered by inclusion) onto itself is strictly decreasing.
\item
  Let E be an ordered set. For each \(x \in E\) let \(U_x\) be the set of all \(y \in E\) such that \(y \geqslant x\) . The mapping \(x \to U_x\) is a strictly decreasing mapping of E into \(\mathfrak{P}(E)\) (ordered by inclusion); indeed, the relation \(x \leqslant y\) is equivalent to \(U_x \supset U_y\) .
\end{enumerate}

An injective monotone mapping of an ordered set E into an ordered set F is strictly monotone; the converse is usually not true, because it may happen that \(f(x) = f(y)\) when neither of the relations \(x \leqslant y\) , \(x \geqslant y\) is true (cf.~no. 12, Proposition 11).

A bijective mapping \(f\) of an ordered set \(\mathbf{E}\) onto an ordered set \(\mathbf{E}'\) is an isomorphism of \(\mathbf{E}\) onto \(\mathbf{E}'\) (no. 3) if and only if both \(f\) and its inverse mapping are increasing.

If I is an ordered index set, a family of subsets \((\mathbf{X}_t)_{i \in I}\) of a set E is said to be increasing if \(\iota \to \mathbf{X}_t\) is an increasing mapping of I into \(\mathfrak{P}(\mathbf{E})\) , ordered by inclusion (in other words, if \(\iota \leqslant x\) implies \(\mathbf{X}_t \subset \mathbf{X}_x\) ). Similarly we define a decreasing, strictly increasing, or strictly decreasing family of subsets \((\mathbf{X}_t)_{i \in I}\) .

PROPOSITION 2. Let E, E' be two ordered sets, and let u: E → E' and v: E' → E be two decreasing mappings such that for all x ∈ E and all x' ∈ E' we have v(u(x)) ≥ x and u(v(x')) ≥ x'. Then

\[
u \circ v \circ u = u \quad \text{and} \quad v \circ u \circ v = v.
\]

For the relation \(v(u(x)) \geqslant x\) implies \((u(v(u(x))) \leqslant u(x)\) because \(u\) is decreasing; on the other hand, we have \(u(v(u(x))) \geqslant u(x)\) by replacing \(x'\) by \(u(x)\) in the inequality \(u(v(x')) \geqslant x'\) . Hence the first equality; the proof of the second is similar.

